\let\negmedspace\undefined
\let\negthickspace\undefined
\documentclass[journal]{IEEEtran}
\usepackage[a5paper, margin=10mm, onecolumn]{geometry}
%\usepackage{lmodern} % Ensure lmodern is loaded for pdflatex
\usepackage{tfrupee} % Include tfrupee package

\setlength{\headheight}{1cm} % Set the height of the header box
\setlength{\headsep}{0mm}     % Set the distance between the header box and the top of the text

\usepackage{gvv-book}
\usepackage{gvv}
\usepackage{cite}
\usepackage{physics}
\usepackage{amsmath,amssymb,amsfonts,amsthm}
\usepackage{algorithmic}
\usepackage{graphicx}
\usepackage{textcomp}
\usepackage{xcolor}
\usepackage{txfonts}
\usepackage{listings}
\usepackage{enumitem}
\usepackage{mathtools}
\usepackage{gensymb}
\usepackage{comment}
\usepackage[breaklinks=true]{hyperref}
\usepackage{tkz-euclide} 
\usepackage{listings}
% \usepackage{gvv}                                        
\def\inputGnumericTable{}                                 
\usepackage[latin1]{inputenc}                                
\usepackage{color}                                            
\usepackage{array}                                            
\usepackage{longtable}                                       
\usepackage{calc}                                             
\usepackage{multirow}                                         
\usepackage{hhline}                                           
\usepackage{ifthen}                                           
\usepackage{lscape}
\begin{document}

\bibliographystyle{IEEEtran}

\title{7.4.29}
\author{EE25BTECH11022 - sankeerthan}
% \maketitle
% \newpage
% \bigskip
{\let\newpage\relax\maketitle}

\renewcommand{\thefigure}{\theenumi}
\renewcommand{\thetable}{\theenumi}
\setlength{\intextsep}{10pt} % Space between text and floats


\numberwithin{align}{enumi}
\numberwithin{figure}{enumi}
\renewcommand{\thetable}{\theenumi}

\textbf{Question:}If one of the diameters of the circle $x^2 + y^2 - 2x - 6y + 6 = 0$ is a chord to the circle with the center $\brak{2,1}$,then the radius of the circle is 
\begin{enumerate}
     \item $\sqrt{3}$
     \item $\sqrt{2}$
     \item $3$
     \item $2$
\end{enumerate}

\textbf{Solution:}
Given the circle equation:
\begin{align}
x^2 + y^2 - 2x - 6y + 6 = 0 \\
\mathbf{x}^T \mathbf{x} + 2\mathbf{u}^T \mathbf{x} + f = 0,\text{where}
\mathbf{x} = \myvec{x \\ y}, \quad \mathbf{u} = \myvec{g \\ f}
\end{align}
From the given equation:
\begin{align}
2\mathbf{u}^T \mathbf{x} = -2x - 6y \implies \mathbf{u} = \myvec{-1 \\ -3}, \quad f = 6
\end{align}
Center of the circle is:
\begin{align}
\mathbf{c_1} = -\mathbf{u} = \myvec{1 \\ 3}
\end{align}
Radius of the circle is:
\begin{align}
r_1 =\sqrt{\mathbf{u}^T \mathbf{u} - f} =\sqrt{(-1)^2 + (-3)^2 - 6} =\sqrt{4} = 2
\end{align}
Given the second circle has center:
\begin{align}
\mathbf{c_2} = \myvec{2 \\ 1}
\end{align}
Assuming the diameter chord line of the first circle with equation $y=3$,
\begin{align}
\mathbf{n} = \myvec{0 \\ 1}, \quad c = 3
\end{align}
Distance from second circle center \(\mathbf{c_2}\) to this chord:
\begin{align}
\frac{|\mathbf{n}^T \mathbf{c_2} - c|}{\|\mathbf{n}\|} = \frac{|1 - 3|}{1} = 2
\end{align}
Therefore, the radius of the second circle  r=2.

\begin{figure}
    \centering
    \includegraphics[width=1\columnwidth]{figs/points.png}
    \caption{}
\end{figure}
\end{document}
